% !TEX root = ../main.tex
\chapter{Numerical Gradient Checking}
\label{ch:gradcheck}

A bug in a custom backward (\chref{ch:custom_functions}) or in a
freshly added primitive's \code{BackwardSpec}
(\chref{ch:op_registry}) is one of the easiest classes of bug to
introduce and one of the hardest to detect by inspection. The
gradient is silently wrong; training appears to proceed; the
model converges to the wrong place or fails to converge at all;
the cause is far upstream of the symptom.

\code{gradcheck} is the standard remedy: compare the analytical
backward against a finite-difference approximation and assert
agreement to within a tolerance. The technique is old (it predates
modern AD by decades) and the implementation is short, but the
practical use requires understanding several subtleties.

This chapter describes \Lucid's \code{lucid.autograd.gradcheck},
its sister \code{gradgradcheck} for second-order verification, and
the choices that go into setting tolerances.

\section{The Idea}
\label{sec:gradcheck:idea}

For a function $f : \mathbb{R}^N \to \mathbb{R}$, the analytical
gradient is $\nabla f$. The two-sided finite-difference
approximation is
\[
\widehat{\partial f}_i = \frac{f(x + h e_i) - f(x - h e_i)}{2h},
\]
where $e_i$ is the $i$-th unit vector and $h$ is a small step
size. \code{gradcheck} computes both and asserts they agree to
within absolute and relative tolerances.

For a vector-output function $f : \mathbb{R}^N \to \mathbb{R}^M$,
the analytical Jacobian is computed via $M$ reverse-mode VJP calls
(one per output, with a one-hot cotangent), and the finite-
difference Jacobian is computed via $2N$ forward calls.
\code{gradcheck} compares the two $M \times N$ matrices entry by
entry.

\subsection{Why the comparison is informative}
\label{ssec:gradcheck:informative}

The two computations use entirely different mechanisms:

\begin{itemize}
  \item Analytical: walks the autograd graph backward, calling
    each op's VJP.
  \item Numerical: calls the forward function $2N$ times with
    perturbed inputs.
\end{itemize}

A bug in either is unlikely to be a bug in both. If the two agree
to $10^{-6}$ relative error, the analytical backward is almost
certainly correct.

\section{The API}
\label{sec:gradcheck:api}

\code{lucid.autograd.gradcheck} signature:

\begin{lstlisting}[style=lucid,language=LucidPython]
def gradcheck(
    func,            # the function to check
    inputs,          # tuple of tensors
    *,
    eps=1e-6,        # finite-difference step size
    atol=1e-5,       # absolute tolerance
    rtol=1e-3,       # relative tolerance
    raise_exception=True,
) -> bool:
    ...
\end{lstlisting}

Typical use:

\begin{lstlisting}[style=lucid,language=LucidPython]
import lucid
from lucid.autograd import gradcheck

x = lucid.randn(5, requires_grad=True, dtype=lucid.float64)
y = lucid.randn(5, requires_grad=True, dtype=lucid.float64)

def f(x, y):
    return (x * y).sum()

assert gradcheck(f, (x, y))
\end{lstlisting}

The call returns True if every gradient entry agrees within
tolerance and raises (or returns False, with
\code{raise\_exception=False}) otherwise.

\subsection{The FP64 convention}
\label{ssec:gradcheck:fp64_conv}

The user passes FP64 tensors, not FP32. The reasoning is the
numerical-error trade-off discussed in
\chref{ch:custom_functions}~Section~\ref{ssec:custom_functions:gradcheck_fp64}:
finite differences with a step of $10^{-6}$ in FP32 lose all
useful precision; in FP64 they retain $\sim 8$ digits.

\Lucid's \code{gradcheck} raises with a helpful message if any
input is FP32 (\code{LucidError: gradcheck requires FP64 inputs
for meaningful tolerance}), suggesting the user cast first.

\section{Tolerance Setting}
\label{sec:gradcheck:tolerance}

The tolerance is the single most important parameter. Too tight,
and a correct backward fails due to the finite-difference error.
Too loose, and a buggy backward passes.

\subsection{Theoretical tolerance}
\label{ssec:gradcheck:theoretical_tol}

For a step size $h$, the central-difference approximation
\[
\widehat{f'}(x) \;=\; \frac{f(x+h) - f(x-h)}{2h}
\]
has total error
\[
\bigl|\widehat{f'}(x) - f'(x)\bigr|
\;\leq\;
\underbrace{\frac{h^2}{6} \bigl|f'''(\xi)\bigr|}_{\text{truncation}}
\;+\;
\underbrace{\frac{\epsilon}{h} \bigl|f(x)\bigr|}_{\text{round-off}}.
\]
The truncation term grows with $h^2$; the round-off term grows
as $1/h$. Their sum has a unique minimum, attained at
\[
h^{*} \;=\; \left(\frac{3\,\epsilon\,|f(x)|}{|f'''(\xi)|}\right)^{1/3}
\;\approx\; \epsilon^{1/3},
\]
giving an optimal error of order $\epsilon^{2/3}$. In FP64
($\epsilon \approx 10^{-16}$), the optimum is $h^* \approx
10^{-5}$ and the floor is $\approx 10^{-10}$. In FP32
($\epsilon \approx 10^{-7}$), the optimum is $h^* \approx
10^{-2}$ and the floor is $\approx 10^{-4}$ --- barely enough
precision to distinguish a correct gradient from a buggy one,
which is why gradcheck insists on FP64.

\Lucid's default $h = 10^{-6}$ is slightly above the optimum; the
default tolerance $\text{atol} = 10^{-5}$ leaves room for a
generous safety margin.

\subsection{Per-op tolerance adjustments}
\label{ssec:gradcheck:per_op_tol}

Some ops have intrinsically larger finite-difference error:

\begin{itemize}
  \item Transcendental ops (\code{exp}, \code{log}) lose
    precision through their internal computation; tolerance must
    be looser, typically \code{atol=1e-4}.
  \item Operations with extreme value range (\code{erfinv} near
    $|x|=1$, \code{logsumexp} with extreme logits) need
    case-specific tolerances; sometimes the function must be
    restricted to a safe input range for gradcheck.
  \item Discontinuous ops (\code{relu} at 0, \code{abs} at 0)
    cannot be gradchecked at the discontinuity; the inputs must
    be chosen to avoid the discontinuous point.
\end{itemize}

\Lucid's test suite has per-op tolerance overrides where needed,
documented in the test files themselves.

\section{What gradcheck Catches}
\label{sec:gradcheck:catches}

A bug in the backward almost always produces a measurable
disagreement. The common failure modes:

\subsection{Sign error}
\label{ssec:gradcheck:sign_error}

A backward with the wrong sign (e.g., \code{return -grad\_out}
instead of \code{return grad\_out}) disagrees by a factor of
exactly 2 (since the relative error is 200\,\%). \code{gradcheck}
catches this immediately.

\subsection{Missing factor}
\label{ssec:gradcheck:missing_factor}

A backward that omits a constant factor (e.g., the $2$ in
$d/dx\, x^2 = 2x$) disagrees by that factor. \code{gradcheck}
catches this immediately.

\subsection{Wrong saved tensor}
\label{ssec:gradcheck:wrong_saved}

A backward that uses the wrong saved quantity (e.g., the input
when it should have been the output) produces a disagreement that
depends on the input range. \code{gradcheck} catches it for
non-pathological inputs.

\subsection{Unbroadcast error}
\label{ssec:gradcheck:unbroadcast_error}

A backward that forgets to sum over the broadcast dimensions (the
unbroadcast step of \chref{ch:binary_functions}) produces a
gradient with the wrong shape, which \code{gradcheck} reports as
a shape mismatch.

\section{What gradcheck Does Not Catch}
\label{sec:gradcheck:misses}

A short list of cases where \code{gradcheck} can pass on a buggy
backward.

\subsection{Errors that are below tolerance}
\label{ssec:gradcheck:below_tol}

A backward that is wrong by less than the tolerance passes. The
remedy is a tighter tolerance, which requires FP64 inputs (which
gradcheck enforces).

\subsection{Errors at points not tested}
\label{ssec:gradcheck:not_tested}

\code{gradcheck} tests at the specific input the user passes. A
bug that manifests only at extreme inputs (very large, very
small, NaN, Inf) may not surface in a randomised test. The
mitigation is to test multiple input regimes (small, moderate,
large) and to spot-check extreme cases manually.

\subsection{Errors in saved-output backward when the output is constant}
\label{ssec:gradcheck:const_out}

If a backward uses a saved output and the test inputs happen to
produce a constant output (or near-constant), the backward may
appear correct even if the output-dependence is wrong. The
mitigation is varied test inputs.

\subsection{Non-deterministic ops}
\label{ssec:gradcheck:nondet}

Operations with non-deterministic backward (\code{scatter\_add}
with duplicate indices on GPU) produce different cotangents on
different runs. \code{gradcheck} cannot pass deterministically for
these; the test must be skipped or run on a deterministic stream.

\section{gradgradcheck}
\label{sec:gradcheck:gradgradcheck}

For higher-order verification, \code{gradgradcheck} checks the
gradient of the gradient. It's the natural extension when the
backward itself contains another backward (e.g., when the model
uses a meta-learning algorithm that requires second-order
gradients).

The mechanism is the same: build the second-order graph via
\code{backward(create\_graph=True)}, then run \code{gradcheck} on
the resulting gradient function.

\subsection{When to gradgradcheck}
\label{ssec:gradcheck:when_gg}

\begin{itemize}
  \item When using \code{hessian}, \code{jacrev o jacrev}, or
    other compositions that require second-order autograd.
  \item When implementing a custom backward that itself uses
    \Lucid\ primitives whose second derivatives must be correct.
\end{itemize}

For ordinary first-order training, gradgradcheck is unnecessary.

\section{The Parity-Test Policy}
\label{sec:gradcheck:parity_policy}

\Lucid's testing policy
(\chref{ch:testing}) requires every public op to have:

\begin{enumerate}
  \item A unit test exercising correctness against a known
    expected output.
  \item A parity test comparing against the reference framework's
    output to within a tight tolerance.
  \item A \code{gradcheck} verifying the backward.
  \item For ops with non-trivial second derivatives, a
    \code{gradgradcheck}.
\end{enumerate}

The discipline catches the easy classes of bug at PR time, before
the op ships. It does not catch the hard cases (numerical
pathologies at extreme inputs, non-determinism), which require
case-by-case attention.

\section{Implementation}
\label{sec:gradcheck:implementation}

\code{gradcheck} is implemented in \code{lucid/autograd/gradcheck.py}.
The core loop:

\begin{lstlisting}[style=lucid,language=LucidPython]
def gradcheck(func, inputs, eps=1e-6, atol=1e-5, rtol=1e-3,
              raise_exception=True):
    # Validate FP64
    for t in inputs:
        if t.dtype != lucid.float64:
            raise LucidError("gradcheck requires float64 inputs")

    # Compute analytical gradients (one per output entry)
    analytical = _analytical_jacobian(func, inputs)

    # Compute numerical gradients via central difference
    numerical = _numerical_jacobian(func, inputs, eps)

    # Compare entry by entry
    for a, n in zip(analytical, numerical):
        diff = (a - n).abs()
        tol = atol + rtol * n.abs()
        if (diff > tol).any():
            if raise_exception:
                raise LucidError(
                    f"gradcheck failed: max diff {diff.max()} "
                    f"exceeds tolerance {tol.max()}")
            return False
    return True
\end{lstlisting}

The full implementation handles edge cases: complex-valued
functions, sparse outputs, the \code{check\_undefined\_grad}
option, and a few others. The core idea is the central-difference
loop above.

\subsection{Cost}
\label{ssec:gradcheck:cost}

A \code{gradcheck} on a function with $N$ inputs and $M$ outputs
costs $2N$ forward calls plus $M$ backward calls. For typical
unit tests with small $N$ and $M$ (single-digit), the cost is
under a second.

For larger inputs ($N$ in the hundreds), the cost is significant
but acceptable for a one-time correctness check. \code{gradcheck}
is not run during normal training; it's a development-time tool.

\section{Worked Example: Validating a Custom Function}
\label{sec:gradcheck:worked}

The custom \code{SwishBeta} from
\chref{ch:custom_functions}~Section~\ref{sec:custom_functions:custom_activation}:

\begin{lstlisting}[style=lucid,language=LucidPython]
import lucid
from lucid.autograd import gradcheck

x = lucid.randn(5, requires_grad=True, dtype=lucid.float64)
beta = lucid.tensor(0.7, requires_grad=True, dtype=lucid.float64)

# Validate first-order gradient
assert gradcheck(SwishBeta.apply, (x, beta), eps=1e-6, atol=1e-5)
print("SwishBeta first-order gradient verified")
\end{lstlisting}

If we had introduced the sign error (\code{return -grad\_x,
-grad\_beta} instead of the correct signs), the test would have
failed with a relative error of approximately 2.0:

\begin{verbatim}
LucidError: gradcheck failed: max diff 1.84e+00 exceeds
tolerance 6.21e-05
\end{verbatim}

The failure points at the input and the offending Function class,
making the bug easy to locate.

\section{Debugging a Failed gradcheck}
\label{sec:gradcheck:debugging_failure}

\code{gradcheck} produces an error message with the maximum
disagreement and its location. The message tells the user what
failed; it doesn't tell why. A short debugging recipe.

\subsection{Step 1: print both gradients}
\label{ssec:gradcheck:debug_print}

Replace \code{gradcheck} with manual computation:

\begin{lstlisting}[style=lucid,language=LucidPython]
import lucid

# Analytical
x = lucid.randn(3, requires_grad=True, dtype=lucid.float64)
y = my_func(x).sum()
y.backward()
print("analytical:", x.grad)

# Numerical (one-sided for brevity)
def num_grad(f, x, eps=1e-6):
    g = lucid.zeros_like(x)
    for i in range(x.numel()):
        x_perturbed = x.clone()
        x_perturbed.flat[i] += eps
        f_plus = f(x_perturbed)
        x_perturbed.flat[i] -= 2*eps
        f_minus = f(x_perturbed)
        g.flat[i] = (f_plus - f_minus) / (2*eps)
    return g

print("numerical:", num_grad(lambda x_: my_func(x_).sum(), x))
\end{lstlisting}

The visual comparison often makes the bug obvious (a sign error
flips, a missing factor shrinks, a wrong saved tensor adds noise
to specific entries).

\subsection{Step 2: simplify the function}
\label{ssec:gradcheck:debug_simplify}

If the function is composed of multiple parts, isolate each. Run
\code{gradcheck} on each part separately. The failing part is
the source of the bug.

\subsection{Step 3: check the saved tensors}
\label{ssec:gradcheck:debug_saved}

A common bug class: backward uses the wrong saved tensor. Print
the saved tensors at the start of \code{backward}, compare to the
forward inputs / outputs you expected. Mismatches surface here.

\subsection{Step 4: check shape}
\label{ssec:gradcheck:debug_shape}

A backward that produces the wrong gradient shape (a missing
\code{sum} after broadcasting, an off-by-one slice) is caught
early by gradcheck's shape check. The error message names the
shapes; line them up against the input shape and the expected
gradient shape.

\section{gradcheck vs Other Verification Tools}
\label{sec:gradcheck:vs_other}

\code{gradcheck} is one tool in a suite of verification
mechanisms.

\subsection{vs parity testing}
\label{ssec:gradcheck:vs_parity}

A parity test (\chref{ch:testing}) compares against a reference
framework's output. It catches forward bugs and backward bugs
together, but it only works for ops where a reference exists. For
genuinely new ops or custom Functions, gradcheck is the only
option.

The two are complementary: parity tests catch ``the answer is
wrong'' bugs; gradcheck catches ``the answer is right but the
gradient is wrong'' bugs.

\subsection{vs anomaly detection}
\label{ssec:gradcheck:vs_anomaly}

\code{lucid.autograd.set\_detect\_anomaly(True)}
(\chref{ch:autograd_engine}) catches NaN gradients at runtime. It
doesn't catch wrong (but finite) gradients; gradcheck does.

\subsection{vs unit testing}
\label{ssec:gradcheck:vs_unit}

A unit test that compares against a hand-derived expected value
catches the same class of bug as gradcheck for specific inputs.
gradcheck generalises: it works against any input the user
chooses, with the analytical machinery doing the comparison
automatically.

\section{The check\_undefined\_grad Option}
\label{sec:gradcheck:undefined_grad}

\code{gradcheck} has an option \code{check\_undefined\_grad=True}
(default True in \Lucid\ 3.5) that additionally verifies the
function handles the case where the upstream cotangent is
\code{None} (undefined). This case arises when the user's
downstream consumer discards an output (e.g., \code{out = (a, b,
c); use(a); del b, c}); the engine passes \code{None} as the
cotangent for the discarded outputs.

The Function's backward should handle this by returning
\code{None} for the corresponding inputs. The check verifies
that the analytical backward does so consistently.

For most user code the check is irrelevant; for library code
shipping in \Lucid's core, it's essential.

\section{Tolerance Recipes}
\label{sec:gradcheck:tolerance_recipes}

A reference for setting tolerances for common situations.

\begin{itemize}
  \item \textbf{Default FP64}: \code{atol=1e-5}, \code{rtol=1e-3}.
  \item \textbf{Transcendental (exp, log, trig)}: \code{atol=1e-4},
    \code{rtol=1e-3}.
  \item \textbf{Reductions over large tensors}: \code{atol=1e-4},
    \code{rtol=1e-3} (accumulated error grows with size).
  \item \textbf{Complex composite (multiple primitives)}:
    \code{atol=1e-4}, \code{rtol=1e-2}.
  \item \textbf{Erfinv, Bessel, special functions}:
    \code{atol=1e-3}, \code{rtol=1e-2} (intrinsic precision
    limit).
  \item \textbf{FP32 (if FP64 not possible)}: \code{atol=1e-2},
    \code{rtol=1e-1} --- barely useful, only catches gross bugs.
\end{itemize}

The general rule: pick the tightest tolerance that the function's
intrinsic precision permits, and tighten the test as the
implementation matures.

\section{Summary and Forward References}
\label{sec:gradcheck:summary}

\code{lucid.autograd.gradcheck} compares the analytical reverse-
mode gradient against a finite-difference approximation, asserting
agreement to within absolute and relative tolerance.
\code{gradgradcheck} extends to second-order verification.
\Lucid's test suite requires gradcheck for every public op, which
catches the easy classes of backward bug at PR time. The tool is
a development-time correctness check, not a training-time
mechanism.

This chapter closes \partref{part:autograd}. \partref{part:math_subsystems}
treats the mathematical subpackages (linalg, FFT, signal,
special functions, distributions, einops) that build on the
autograd machinery covered here.

\begin{lucidnote}
For the user: \code{gradcheck} is the right tool when you have
written a custom \code{Function} and want to verify its backward.
The standard recipe is FP64 inputs, the default tolerances, and
\code{raise\_exception=True} (the default) to surface failures
loudly. For first-order, \code{gradcheck}; for second-order,
\code{gradgradcheck}.
\end{lucidnote}
